\chapter{A Simulation's Life Cycle}
\label{ch:lifecycle}

\chapterorient{In the last chapter we stood outside the machine and watched a
configuration go in and a physical product come out. Now we open the lid just
far enough to see the stages the run passes through---parse, configure,
dispatch, build the mesh, solve, refine, output---and what travels between them.
One stage will turn out to be the hinge on which all six analyses swing, and one
will be our first glimpse of the single most important structure in the book:
the \emph{fold}. By the end you will be able to trace any Palace run from the
press of \textsf{enter} to the file it writes, and to point to the exact line
where ``one program'' becomes ``six.''}

\section{From black box to pipeline}

\Cref{ch:targets} showed every analysis as one black box: configuration in,
product out. That picture is true but coarse. If we slice the box open along the
direction the data flows, we find not one step but a short \emph{pipeline} of
them, run in sequence, the same sequence every time, for every analysis. This
chapter is a tour of that pipeline---the \emph{lifecycle} of a simulation.

The lifecycle is worth a chapter of its own for two reasons. First, it is the
\emph{shared skeleton}: all six analyses pass through exactly these stages in
exactly this order, so understanding the lifecycle once means understanding the
top-level shape of all six at once. Second, two of its stages carry the book's
central ideas in seed form---the \emph{dispatch} that turns one program into six,
and the \emph{adaptive loop} that gives us our first fold. We meet both here, in
their natural habitat, before we study either in depth.

In real source, the lifecycle lives in two files. A short driver, \code{main.cpp},
sets the run up and chooses which analysis to perform; a shared base class,
\code{basesolver.cpp}, runs the chosen analysis under an adaptive outer loop.
We will not read C++ in this book---we read the lifecycle as a sequence of pure
transformations on values---but it helps to know that the two halves we are
about to describe are two real, separate pieces of the program.

\section{The six stages}

\Cref{fig:lifecycle} is the whole chapter in one picture. A run proceeds left to
right through six stages. Most of the arrows are ordinary forward flow; one is
dashed and points \emph{backward}---the adaptive loop, the only place the
pipeline ever revisits a stage. Keep the figure in view; the subsections below
walk the stages one at a time.

\begin{figure}[t]
\centering
\begin{tikzpicture}[node distance=7mm and 9mm]
  \node[stage] (parse) {parse\\config};
  \node[stage, right=of parse] (dev) {configure\\device};
  \node[stage, right=of dev] (disp) {dispatch on\\problem type};
  \node[stage, right=of disp] (mesh) {build\\mesh};
  % the solve / estimate / refine band
  \node[opbox, below=12mm of disp, xshift=10mm] (solve) {Solve};
  \node[opbox, right=of solve] (est) {estimate};
  \node[opbox, right=of est] (mark) {mark};
  \node[opbox, right=of mark] (ref) {refine};
  \node[product, below=12mm of mesh, xshift=22mm] (out) {output\\product};
  % top row forward flow
  \draw[flow] (parse) -- (dev);
  \draw[flow] (dev) -- (disp);
  \draw[flow] (disp) -- (mesh);
  % into the solve band
  \draw[flow] (mesh.south) |- (solve.west);
  % the adaptive band, left to right
  \draw[flow] (solve) -- node[above,font=\scriptsize]{soln.} (est);
  \draw[flow] (est) -- node[above,font=\scriptsize]{$\eta_K$} (mark);
  \draw[flow] (mark) -- (ref);
  % the back-edge: refine -> solve (the AMR loop)
  \draw[backflow] (ref.south) -- ++(0,-6mm) -| node[pos=0.25,below,font=\scriptsize]{refined mesh}
    (solve.south);
  % exit to output
  \draw[flow] (est.south) -- ++(0,-3mm) -| (out.north);
  % a faint band behind the loop
  \begin{scope}[on background layer]
    \node[band, fill=simBgGold, fit=(solve)(est)(mark)(ref),
      inner sep=6pt] (loopband) {};
  \end{scope}
  \node[cornertag, above=0.5pt of loopband.north] {the adaptive estimate--mark--refine loop};
\end{tikzpicture}
\caption[The simulation lifecycle]{The lifecycle of a single run. Four set-up
stages along the top---\emph{parse}, \emph{configure}, \emph{dispatch},
\emph{build mesh}---feed into the shaded \emph{solve--estimate--mark--refine}
loop, whose dashed back-edge is the only place the pipeline revisits a stage.
When adaptive refinement is switched off, the loop runs its body exactly once
and the back-edge is never taken; what remains is a straight line from config to
product. The hinge stage is \emph{dispatch}, expanded in \cref{fig:dispatch}.}
\label{fig:lifecycle}
\end{figure}

\subsection{Parse the configuration}

The run begins with a single file: a configuration written by the user. Parsing
it produces one in-memory value---call it the \code{config}---that holds
everything the run needs to know: the mesh file and polynomial order, the
material properties, the boundary excitations, the refinement settings, and,
crucially, the \emph{problem type} naming which of the six analyses to perform.
In Palace this value is an \code{IoData} record built in one line near the top of
\code{main.cpp}.

The single most important fact about the \code{config} is that it is
\emph{read-only}. It is built once, at the start, and then \emph{never modified}.
Every later stage consults it but none writes to it. We will lean on this fact
repeatedly: because the configuration cannot change mid-run, every stage that
depends only on the configuration computes the same thing no matter when we run
it, which is exactly the property that lets us reorder, hoist, and parallelize
the later stages with confidence.

\begin{notationbox}
We use \code{config} for the parsed configuration value throughout the book. It
is a \emph{record}---a bundle of named fields, like
\lstinline[style=l4style]|{problem_type, mesh, materials}|---and we read its
fields with a dot, as in \lstinline[style=l4style]|config.problem_type|. Records
get their full treatment in \cref{ch:bigidea}; for now, think of \code{config} as
a read-only struct.
\end{notationbox}

\subsection{Configure the device}

Before any arithmetic happens, the run brings up the hardware and numerical
libraries it will compute on: it selects a \emph{device} (a CPU or a GPU),
initializes the tensor backend that will run the element kernels, and starts the
linear-algebra and eigenvalue libraries. All of these choices are read from the
\code{config}; none of them touches the physics.

This stage is mostly plumbing, and we will say little more about it---but its
\emph{existence} is the seed of one of the book's themes. The reason Palace can
ask, in one line of configuration, to run the very same simulation on a laptop
CPU or a datacenter GPU is that the computation downstream is expressed
\emph{against an abstract device}, not against any particular one. The functional,
immutable-tensor presentation this book develops is precisely the shape that
makes a single description run on many devices. We will return to this point in
\cref{ch:bigidea}; here, simply note that ``configure the device'' is where the
backend is chosen, and that everything after it is device-agnostic.

\subsection{Dispatch on the problem type}

Now comes the hinge. The run reads one field of the \code{config}---the problem
type---and on the strength of that single value selects which of the six
analyses to perform. In source it is a \code{switch} on
\lstinline[style=l4style]|config.problem_type| with six branches; in our reading
it is a function that maps a problem type to a \emph{driver}, the object that
knows how to run one analysis end to end:
\begin{lstlisting}[style=l4style]
dispatch :: ProblemType -> Driver
dispatch ELECTROSTATIC = electrostaticDriver
dispatch MAGNETOSTATIC = magnetostaticDriver
dispatch EIGENMODE     = eigenmodeDriver
dispatch DRIVEN        = drivenDriver
dispatch TRANSIENT     = transientDriver
dispatch BOUNDARYMODE  = boundaryModeDriver
\end{lstlisting}

This is the place where ``one program'' becomes ``six.'' Everything before
dispatch is shared by all analyses; everything after is, in part, specialized to
the one we chose. We call this point the \emph{specialization seam}, and it is so
central that it gets its own figure (\cref{fig:dispatch}) and its own section
(\cref{sec:seam}) below.

\begin{figure}[t]
\centering
\begin{tikzpicture}[node distance=4mm]
  \node[stage, minimum width=24mm] (sw) {\code{switch}\\\footnotesize\itshape problem type};
  % six driver boxes in a fan
  \node[opbox, above right=10mm and 22mm of sw] (es) {electrostatic};
  \node[opbox, above right=3mm  and 22mm of sw] (ms) {magnetostatic};
  \node[opbox, right=22mm of sw, yshift=2mm]    (em) {eigenmode};
  \node[opbox, below right=3mm  and 22mm of sw] (dr) {driven};
  \node[opbox, below right=10mm and 22mm of sw] (tr) {transient};
  \node[opbox, below right=17mm and 22mm of sw] (bm) {boundary-mode};
  \draw[flow] (sw.east) -- (es.west);
  \draw[flow] (sw.east) -- (ms.west);
  \draw[flow] (sw.east) -- (em.west);
  \draw[flow] (sw.east) -- (dr.west);
  \draw[flow] (sw.east) -- (tr.west);
  \draw[flow] (sw.east) -- (bm.west);
  % one-word hints
  \node[cornertag, right=2mm of es] {capacitance};
  \node[cornertag, right=2mm of ms] {inductance};
  \node[cornertag, right=2mm of em] {resonance};
  \node[cornertag, right=2mm of dr] {S-parameters};
  \node[cornertag, right=2mm of tr] {time march};
  \node[cornertag, right=2mm of bm] {waveguide};
\end{tikzpicture}
\caption[The dispatch fan-out]{The specialization seam. A single
\code{switch} on the problem type fans out to six drivers, one per analysis,
each tagged with the product it computes. This is the only branch in the entire
lifecycle that depends on \emph{which} analysis we are running; every other stage
is shared. The six drivers are the subject of \cref{part:modalities}; that they
share one skeleton is the claim \cref{ch:situating} makes precise.}
\label{fig:dispatch}
\end{figure}

\subsection{Build the mesh}

With a driver chosen, the run builds the computational mesh: it loads the
geometry file named in the \code{config}, applies any driver-specific
preprocessing, partitions it, and performs any \emph{a-priori} refinement the
user requested. The output is a \code{mesh}---the discretized geometry on which
the fields will live.

Two things about this stage matter for the chapters ahead. First, the mesh is
the first value the lifecycle \emph{produces} (as opposed to merely reads): the
\code{config} came from the user, but the \code{mesh} is computed. Second, mesh
construction is \emph{driver-agnostic in shape}: every analysis builds a mesh the
same way, even though the function spaces later laid over it differ. We treat the
mesh as the second of the values that flow between stages, after the config.

\subsection{Solve}

Here the actual physics happens. The chosen driver takes the mesh and the
configuration, assembles the operators the analysis needs, solves the resulting
system, and reduces the solution to the analysis's physical product. This is the
\emph{assemble--solve--reduce} skeleton we met in \cref{ch:targets}, now occupying
a single stage of the lifecycle. Its internals are different for each driver---an
electrostatic solve is a handful of linear solves, an eigenmode solve is one
opaque eigenvalue call, a transient solve is a march through time---and unpacking
those differences is the work of \cref{part:modalities}.

For the lifecycle, the solve stage produces two outputs, and it is worth being
precise about both. The obvious one is the \emph{product}: the capacitance
matrix, the S-parameters, the field history---the thing the user ran the
simulation to obtain. The less obvious one is a set of \emph{error
indicators}: one nonnegative number $\eta_K$ per mesh element, estimating how
much of the solution's error is concentrated there. The product is what the user
wants; the error indicators are what the \emph{next} stage wants.

\subsection{Estimate, mark, refine---the adaptive loop}
\label{sec:amrloop}

The final structural stage is the one that makes \cref{fig:lifecycle} a loop
rather than a line. Having solved once and obtained error indicators, the run
asks a question: is the solution accurate enough? If the total error is below the
user's tolerance, it stops and outputs the product. If not, it improves the mesh
where the error is worst and solves again. Concretely:

\begin{enumerate}[nosep]
  \item \textbf{Estimate.} Turn the solution into the per-element error
  indicators $\eta_K$ (this is folded into the solve stage above; we name it
  separately because it is conceptually its own step).
  \item \textbf{Mark.} Select the subset of elements carrying the bulk of the
  error---a small set responsible for a large fraction of $\sum_K \eta_K^2$.
  \item \textbf{Refine.} Subdivide the marked elements, producing a finer mesh
  concentrated exactly where the solution needs the most resolution.
  \item \textbf{Repeat.} Solve again on the refined mesh, re-estimate, and loop,
  until the error falls below tolerance or a resource budget is exhausted.
\end{enumerate}

This is \emph{adaptive mesh refinement} (AMR), and it is studied in full in
\cref{ch:amr}. What concerns us here is its \emph{shape}, because that shape is
the most important structure in this book.

\subsection{Finalize and output}

When the loop terminates, the run writes the product to disk, records run
metadata (timings, memory, the configuration used), and tears down the device
and libraries. The product itself---the capacitance, the modes, the field
history---is written by the driver, not by the top-level \code{main}; \code{main}
writes only the metadata. This matters less for the physics than for the
bookkeeping: it tells us that the ``output'' of the lifecycle is exactly the
driver's product, and that the lifecycle ROOT contributes the \emph{scaffold}
around it, not the product itself.

\section{What flows between the stages}

We have named the stages; now read the figure the other way and ask what travels
the arrows. Four kinds of value flow through the lifecycle, and keeping them
straight is half of understanding it.

\begin{description}[style=nextline, leftmargin=0pt, font=\normalfont\itshape]
  \item[The configuration] enters at parse and is \emph{read-only} thereafter.
  Every stage may consult it; none may change it. It is the one value present
  from the first stage to the last.
  \item[The mesh] is produced by the build stage and consumed by every solve.
  It is the one value the \emph{loop} modifies: each refinement replaces the mesh
  with a finer one. Outside the loop it is built once and read.
  \item[The solution and its error indicators] are produced by each solve. The
  solution becomes (after reduction) the product; the error indicators drive the
  mark--refine decision. They are the values that \emph{close the loop}---the
  output of one iteration that becomes (via the refined mesh) the input of the
  next.
  \item[The product] is the final output: the physical quantity the user asked
  for. It is produced by the driver and written out at finalize.
\end{description}

\begin{keyidea}
The configuration is read-only and lives for the whole run; the mesh is built
once and then \emph{refined in the loop}; the solution and error indicators are
recomputed every iteration; the product is written at the end. The only value
that genuinely \emph{threads through} the adaptive loop---changing from one
iteration to the next---is the mesh (carrying its indicators). Everything else is
either fixed for the run or recomputed from scratch each pass. That single
threaded value is what makes the loop a \emph{fold}.
\end{keyidea}

\section{The specialization seam: one program, six analyses}
\label{sec:seam}

We promised in \cref{ch:targets} that the six analyses are one machine with parts
swapped. The lifecycle lets us say exactly \emph{where} the swap happens: at the
dispatch stage, and almost nowhere else.

Read \cref{fig:lifecycle} again with this question in mind: which stages depend on
\emph{which} analysis we are running? Parsing the config does not---it reads the
same file format for all six. Configuring the device does not. Building the mesh
does not, in shape. The adaptive loop does not---estimate, mark, and refine are
the same procedure regardless of what was solved. Even \emph{finalize} is
uniform. The \emph{only} stage that branches on the problem type is dispatch, and
the only stage whose internals differ between analyses is the solve it selects.

This is why we call dispatch the \emph{specialization seam}. Above it, the program
is one shared skeleton; below it, a driver fills in the one stage---solve---that
genuinely differs. The drivers are not six unrelated programs sharing a launcher;
they are six fillings of one mould. \Cref{ch:situating} makes this precise by
showing that even the solve stage varies in a tightly constrained
way---there are only \emph{four kinds} of solve among the six analyses---and that
the reductions that follow come in only three shapes.

\begin{pitfall}
It is tempting to read ``six drivers'' as ``six programs'' and conclude that the
shared lifecycle is just a convenience---a common \code{main} that happens to
launch any of them. That gets the emphasis backwards. The lifecycle is not
incidental scaffolding around six independent solvers; it is the \emph{primary
structure}, and the drivers are specializations of its one variable stage. The
whole book is organized around this inversion: we study the shared skeleton first
(in \cref{ch:situating} and the parts that follow) and the per-analysis
differences last (\cref{part:modalities}), because the shared part is larger and
more fundamental than the differences.
\end{pitfall}

\section{First sight of the fold}
\label{sec:fold}

Return to the adaptive loop of \cref{sec:amrloop}, and look past what it
\emph{does} to the \emph{shape} of how it does it. We have a starting value (the
initial mesh). We repeatedly apply a step that takes the current value and
produces the next one (solve, estimate, mark, refine $\to$ a finer mesh). We stop
when a condition on the accumulated state holds (the error is small enough, or
the budget is spent). And we report a result derived along the way (the product
of the final solve).

A starting value, a step that carries state forward, a stopping condition, a
result threaded through: this pattern is a \emph{fold}. It is the single most
reused structure in this book, and the adaptive loop is your first encounter with
it. In the framework's vocabulary the combinator is called \code{fold\_solve},
and the adaptive loop is written, schematically,
\begin{lstlisting}[style=l4style]
-- run the driver under the adaptive estimate-mark-refine loop
adaptive cfg drv mesh0 = fold_solve step mesh0
  where step m = let (product, etas) = drv cfg m   -- solve + estimate
                     m'               = refine (mark etas) m
                 in  (product, m')                 -- carry the refined mesh forward
\end{lstlisting}

The crucial property of a fold---the one that distinguishes it from a simple
loop over independent items---is that each step \emph{depends on the previous}.
The mesh fed to iteration $n+1$ is exactly the mesh that iteration $n$ produced;
the iterations do not commute and cannot be run in parallel or reordered. This is
in sharp contrast to, say, the electrostatic analysis's inner solves, which run
the \emph{same} operator against several independent right-hand sides and could
be done in any order. The distinction between an independent family of solves (a
\emph{map}) and a state-threaded chain of solves (a \emph{fold}) is one of the
four ``solve corners'' \cref{ch:situating} will name, and the fold is studied as a
combinator in its own right in \cref{ch:fold}.

\begin{keyidea}
The adaptive loop is a \emph{fold}: a step applied repeatedly, each pass
consuming the value the previous pass produced, terminating on a condition over
the accumulated state. The threaded value is the mesh. This same fold structure
will reappear as the engine of time-stepping, of Krylov solvers, and of
eigenvalue iterations---once you can see the fold here, you will see it
everywhere.
\end{keyidea}

\subsection{When the loop degenerates to a line}

There is one more observation, small but clarifying. Adaptive refinement is
\emph{optional}: a user who wants a single solve on a fixed mesh simply sets the
maximum refinement count to zero. When they do, the loop's stopping condition is
true before the first refinement, so the body runs exactly once: solve, estimate,
done. The back-edge in \cref{fig:lifecycle} is never taken, and the lifecycle
collapses to a straight line---config, device, dispatch, mesh, one solve, output.

This is worth dwelling on because it means the \emph{non-adaptive} lifecycle is
not a different program; it is the adaptive lifecycle with its fold run for a
single iteration. A fold of one step is just that step. So the electrostatic
analysis of \cref{ch:targets}, which most readers will picture as a simple
straight-through computation, is in truth the very same adaptive lifecycle---it
has just degenerated to a single pass. We lose nothing by always picturing the
loop; the line is its special case.

\begin{example}
A capacitance extraction with refinement disabled runs: parse the config;
configure the CPU; dispatch to the electrostatic driver; build the mesh; solve
once (yielding the capacitance matrix and a set of error indicators that nobody
will consult); find that the refinement budget is zero, so skip refinement; write
the capacitance matrix; finalize. Six stages, the loop run once. Re-enable
refinement and the only change is that the solve--estimate--mark--refine band
repeats until the indicators fall below tolerance---the same program, its fold
unrolled further.
\end{example}

\section{Where this leaves us}

We have opened the black box of \cref{ch:targets} into a pipeline of six stages
and identified its two load-bearing features: the \emph{dispatch} that
specializes one shared skeleton into six analyses, and the \emph{adaptive loop}
whose shape is a fold. We have catalogued the four kinds of value that flow
between stages, and we have seen that the familiar straight-through simulation is
just the adaptive lifecycle with its fold run once.

Two threads now lead forward. The \emph{specialization} thread asks: given that
the analyses differ only at the solve stage, exactly how do they differ?
\Cref{ch:situating} answers with the four solve corners and three reduction
shapes---the map that the rest of the book fills in. The \emph{representation}
thread asks: what is this ``value-threading,'' ``read-only configuration,''
``fold'' language we have started speaking, and why is it the right one for a
simulator? \Cref{ch:bigidea} answers that, taking up the shift from in-place
array mutation to pure functions over immutable tensors that the whole book is
built on.

\summarysection
Every Palace run follows the same lifecycle: parse the read-only configuration;
configure the compute device; \emph{dispatch} on the problem type to one of six
drivers; build the mesh; run the driver's \emph{solve} (assemble--solve--reduce),
which yields both a product and per-element error indicators; and run the
\emph{adaptive estimate--mark--refine loop}, refining the mesh where the error is
worst and re-solving until the error is small enough, before writing the product.
Four kinds of value flow through: the read-only config (whole run), the mesh
(built once, refined in the loop), the solution and its error indicators
(recomputed each pass), and the product (written at the end). Dispatch is the
\emph{specialization seam}---the one stage that branches on the analysis---so the
six analyses are one shared skeleton with a single swapped stage. The adaptive
loop is a \emph{fold}: a step that threads the mesh forward, our first sight of
the structure that organizes the entire book; when refinement is off, the fold
runs once and the lifecycle is a straight line.

\furtherreading
The lifecycle described here is implemented in the Palace solver's top-level
driver and its shared solver base class~\citep{palace2023}; readers comfortable
with C++ can profitably read \code{main.cpp} (the set-up and dispatch) alongside
this chapter. The adaptive-refinement loop rests on a-posteriori error estimation
and bulk marking, due to Zienkiewicz and Zhu~\citep{zienkiewiczzhu1987} and
D\"orfler~\citep{dorfler1996} respectively; we develop both in \cref{ch:amr}. The
fold, the structural heart of \cref{sec:fold}, is a staple of functional
programming; a reader curious to get ahead may consult any introduction to the
subject~\citep{peytonjones2003}, though we build the idea from nothing in
\cref{ch:fold}.

\exercisessection
\begin{enumerate}[nosep]
  \item List the six stages of the lifecycle in order. For each, state whether
  its behaviour depends on which of the six analyses is being run. (Exactly one
  should depend on it in its \emph{branching}, and one more in its
  \emph{internals}.)
  \item The configuration is read-only for the entire run. Name two concrete
  benefits this gives---one about reasoning, one about hardware. (Hint: re-read
  the parse and configure-device subsections.)
  \item A user sets the maximum refinement count to zero. Trace the lifecycle
  step by step and explain, in terms of the fold, why the result is a single
  straight-through solve. How many times does the step run?
  \item Of the four kinds of value that flow between stages, exactly one is
  threaded through the adaptive loop (changed from one iteration to the next).
  Which? Why does that single threaded value make the loop a fold rather than a
  map?
  \item The dispatch stage is the only one that branches on the problem type.
  Suppose someone proposed adding a seventh analysis to Palace. Using
  \cref{fig:lifecycle}, predict which stages they would have to touch and which
  they could leave untouched. (We test this prediction against reality in
  \cref{part:modalities}.)
\end{enumerate}
